% Standard pen-and-paper worksheet preamble, in the form m01 and m02 use.
% Compile with: xelatex -interaction=nonstopmode exercise.tex   (run twice)
%
% 14pt, not 12pt: the sheet is read with a pencil in hand. There are no blanks
% and no ruled answer gaps anywhere on it -- questions are asked in words and
% answered on the back -- so the paper those would have cost goes to the text
% and to the drawings.
\documentclass[a4paper, 14pt]{extarticle}
\usepackage[top=0.75in, bottom=0.7in, left=0.85in, right=0.85in]{geometry}
\usepackage{amsmath}
\usepackage{amssymb}
\usepackage{array}
\usepackage{graphicx}
\usepackage[hidelinks]{hyperref}   % hidelinks, or the print carries colour boxes
\usepackage{tcolorbox}
\usepackage{enumitem}
\usepackage{etoolbox}
\usepackage{tikz}
\usepackage{fontspec}
\usetikzlibrary{calc,decorations,patterns,arrows,decorations.pathmorphing,positioning}
\definecolor{pltblue}{HTML}{1F77B4}

% Body face: Charter, and the four ways to get it. XCharter *is* Charter and
% ships with a full TeX Live, so it is asked for first and by file -- a CI
% runner answers yes to \IfFontExistsTF{Charter}, then cannot typeset it, which
% is how these sheets stopped building there.
\IfFontExistsTF{XCharter-Roman.otf}{
  \setmainfont{XCharter}[
    Extension      = .otf,
    UprightFont    = *-Roman,
    BoldFont       = *-Bold,
    ItalicFont     = *-Italic,
    BoldItalicFont = *-BoldItalic]
}{
  \IfFontExistsTF{Charter}{\setmainfont{Charter}}{
    \IfFontExistsTF{Iowan Old Style}{\setmainfont{Iowan Old Style}}{
      \IfFontExistsTF{Georgia}{\setmainfont{Georgia}}{}}}}

% Handwriting face, resolved once into \ppfont. Used for a heading or two (the
% answer key's "Answers"), not for the body and not for the drawings: the
% figures are read closely and want the text face.
%
% Naming a font that is not installed is a hard error, not a fallback, which is
% why this is a chain and not a line: seven sheets used to name one outright and
% stopped building the day it was uninstalled. Write \ppfont, never
% \fontspec{<a name>}. Excalifont is vendored in tools/fonts/ and installed by
% CI, so it is the one face that is always present.
\newcommand{\ppfont}{}
\IfFontExistsTF{Pretty Neat}{\renewcommand{\ppfont}{\fontspec{Pretty Neat}}}{
  \IfFontExistsTF{Humor Sans}{\renewcommand{\ppfont}{\fontspec{Humor Sans}}}{
    \IfFontExistsTF{Excalifont}{\renewcommand{\ppfont}{\fontspec{Excalifont}}}{
      \renewcommand{\ppfont}{}}}}

\setlength{\parindent}{0pt}
\setlength{\parskip}{0.38em}
\setlength{\emergencystretch}{1.2em}
\raggedbottom

% Break the page unless #1 of it is left.
%
% This is needspace's glue trick, written out because needspace.sty lives in
% texlive-latex-extra and not in a BasicTeX install. Do NOT rewrite it as a test
% on \pagegoal minus \pagetotal: the page builder has not necessarily run when
% the macro expands, so those numbers can still be the previous page's, the test
% says "no room" on a page that is empty, and the \newpage throws that page
% away. m02 shipped a blank page 3 that way, and it is invisible in the source.
\newcommand{\needroom}[1]{%
  \par
  \begingroup
    \setlength{\dimen0}{#1}%
    \vskip 0pt plus \dimen0
    \penalty -100%
    \vskip 0pt plus -\dimen0
  \endgroup}

% A part heading. The sheet has no title: it opens on Part 1. The heading
% carries the break test, because a heading printed at the foot of a page with
% nothing under it is worse than a short page.
\newcommand{\parthead}[2]{%
  \needroom{6\baselineskip}%
  \partheadhere{#1}{#2}}

% The same heading with no break test, for one already inside a minipage --
% \newpage means nothing in there.
\newcommand{\partheadhere}[2]{%
  \vspace{0.5em}
  {\large\bfseries Part #1 --- #2}\par
  \vspace{-0.4em}\rule{\textwidth}{0.8pt}\par
  \vspace{0.1em}}

% Centred fixed-width cell, for tables the student writes into.
\newcolumntype{C}[1]{>{\centering\arraybackslash}p{#1}}

% The lab the sheet hands off to, if it has one. Both the QR code and the
% printed address are the course's own short link, never the notebook's: paper
% outlives any one upload, so re-uploading a lab is a line on the server
% (/etc/caddy/Caddyfile, block go.skojaku.com) and not a reprint.
%
%   uvx --from segno segno --output=m07lab-qr.png --scale=20 --border=4 \
%       --error=m "https://go.skojaku.com/m07lab"
%
% --border=4 is the quiet zone the QR spec asks for and is not optional: at
% --border=1 some codes run to the edge of the PNG and OpenCV reads nothing.
\newcommand{\molablink}{https://go.skojaku.com/m07lab}
\newcommand{\molaburl}{\href{\molablink}{\texttt{go.skojaku.com/m07lab}}}
% The simulator, in the view the sheet needs (?sheet=1: the degree law hidden,
% the most and the least visited node named). Same Caddy block as the lab.
\newcommand{\simlink}{https://go.skojaku.com/m07sim}
\newcommand{\simurl}{\href{\simlink}{\texttt{go.skojaku.com/m07sim}}}

% xkcd-style hand-drawn line decoration.
% WARNING: applying [xkcd] to many nodes can throw "Dimension too large".
% Use it on a handful of \draw lines only, never on every node.
\makeatletter
\pgfset{
  /pgf/decoration/randomness/.initial=2,
  /pgf/decoration/wavelength/.initial=100
}
\pgfdeclaredecoration{sketch}{init}{
  \state{init}[width=0pt,next state=draw,persistent precomputation={
    \pgfmathsetmacro\pgf@lib@dec@sketch@t0
  }]{}
  \state{draw}[width=\pgfdecorationsegmentlength,
  auto corner on length=\pgfdecorationsegmentlength,
  persistent precomputation={
    \pgfmathsetmacro\pgf@lib@dec@sketch@t{mod(\pgf@lib@dec@sketch@t+pow(\pgfkeysvalueof{/pgf/decoration/randomness},rand),\pgfkeysvalueof{/pgf/decoration/wavelength})}
  }]{
    \pgfmathparse{sin(2*\pgf@lib@dec@sketch@t*pi/\pgfkeysvalueof{/pgf/decoration/wavelength} r)}
    \pgfpathlineto{\pgfqpoint{\pgfdecorationsegmentlength}{\pgfmathresult\pgfdecorationsegmentamplitude}}
  }
  \state{final}{}
}
\tikzset{xkcd/.style={decorate,decoration={sketch,segment length=0.5pt,amplitude=0.5pt}}}
\makeatother
\definecolor{pltrust}{HTML}{B14434}
\usetikzlibrary{calc}
% The three figures the Module 7 sheet draws, kept in their own file because
% the answer key draws exactly the same figures with the answers on them.
% Needs tikz and the colours pltblue and pltrust, which the including file sets
% up.
% ---------------------------------------------------------------------------
\tikzset{
  friend/.style  = {line width=1.3pt, black!58, line cap=round},
  person/.style  = {circle, draw=black!72, fill=white, line width=0.9pt,
                    minimum size=0.74cm, inner sep=0pt},
  % The walker's start: a second ring round the person, as on slide 6.
  start/.style   = {person, double=white, double distance=1.6pt},
  % A box the student writes a number into (a share, in Question 6).
  sbox/.style    = {draw=pltblue!70, line width=0.7pt, fill=white,
                    minimum width=1.25cm, minimum height=0.62cm, inner sep=0pt,
                    rounded corners=1.2pt},
  member/.style  = {circle, draw=black!72, fill=white, line width=0.9pt,
                    minimum size=0.66cm, inner sep=0pt, font=\small},
  dot/.style     = {circle, fill=pltblue!80!black, inner sep=0pt,
                    minimum size=0.22cm},
}

% ---------------------------------------------------------------------------
% Part 1: four friends. A--B, A--C, B--C, A--D. Degrees 3, 2, 2, 1, so 2m = 8
% and the long-run shares are 3/8, 2/8, 2/8, 1/8. D hangs off A, so the walker
% swings between A and the rest; the share column of Question 6 does not care.
%
% #1 is drawn last, on top: the answer key writes the shares into the boxes
% with \node at (bA) {...}; and so on.
\newcommand{\fournet}[1]{%
\begin{tikzpicture}
  \node[start]  (A) at (0, 0)      {A};
  \node[person] (B) at (1.9, 1.25) {B};
  \node[person] (C) at (2.6, -0.45){C};
  \node[person] (D) at (-1.9, 0)   {D};
  \draw[friend] (A) -- (B);
  \draw[friend] (A) -- (C);
  \draw[friend] (B) -- (C);
  \draw[friend] (A) -- (D);
  \node[font=\footnotesize, black!65, above=0.08cm of A] {start};
  \node[sbox] (bA) at (0, -1.05)   {};
  \node[sbox] (bB) at (3.35, 1.25) {};
  \node[sbox] (bC) at (2.6, -1.45) {};
  \node[sbox] (bD) at (-1.9, -1.0) {};
  #1
\end{tikzpicture}}

% ---------------------------------------------------------------------------
% Part 3: twelve members, the network of slides 14-45. Members 0-5 and 6-11
% are two groups of six, everybody friends with everybody in the group except
% 0-3 and 7-10, and two friendships cross: 2-8 and 5-11. 30 friendships.
% The layout is the deck's (slides/m07/figures/smallgraphs.py, TWO12_XY): two
% hexagons, the missing friendships are long diagonals, and the two crossing
% friendships are parallel horizontal lines that touch no other member.
%
% #1 is the style of members 0-5 and #2 of members 6-11 (empty on the sheet:
% the student darkens them; the key fills one side).
\newcommand{\tworooms}[2]{%
\begin{tikzpicture}[x=0.8cm, y=0.8cm]
  \foreach \n/\a in {2/60, 1/120, 0/180, 4/240, 5/300, 3/0}
    \coordinate (m\n) at ($(1.5, 2.0) + (\a:1.7)$);
  \foreach \n/\a in {6/60, 8/120, 10/180, 11/240, 9/300, 7/0}
    \coordinate (m\n) at ($(7.5, 2.0) + (\a:1.7)$);
  \foreach \i/\j in {0/1, 0/2, 0/4, 0/5, 1/2, 1/3, 1/4, 1/5, 2/3, 2/4, 2/5,
                     3/4, 3/5, 4/5,
                     6/7, 6/8, 6/9, 6/10, 6/11, 7/8, 7/9, 7/11, 8/9, 8/10,
                     8/11, 9/10, 9/11, 10/11,
                     2/8, 5/11}
    \draw[friend] (m\i) -- (m\j);
  \foreach \n in {0,...,5}  \node[member, #1] at (m\n) {\n};
  \foreach \n in {6,...,11} \node[member, #2] at (m\n) {\n};
\end{tikzpicture}}

% ---------------------------------------------------------------------------
% Part 3: the empty axes the second eigenvector is plotted on, one column per
% member, -1 to 1. #1 is drawn on top (the key's dots).
\newcommand{\eigaxes}[1]{%
\begin{tikzpicture}[x=1.15cm, y=2.9cm]
  \foreach \y in {-1, -0.5, 0.5, 1}
    \draw[black!14, line width=0.5pt] (-0.5, \y) -- (11.5, \y);
  \foreach \x in {0,...,11}
    \draw[black!14, line width=0.5pt] (\x, -1.12) -- (\x, 1.12);
  \draw[black!70, line width=0.8pt, dashed] (-0.5, 0) -- (11.5, 0);
  \draw[black!80, line width=0.9pt] (-0.5, -1.12) -- (-0.5, 1.12);
  \foreach \y/\l in {-1/{$-$1}, -0.5/{$-$0.5}, 0/{0}, 0.5/{0.5}, 1/{1}}
    \node[left, font=\small] at (-0.5, \y) {\l};
  \foreach \x in {0,...,11}
    \node[anchor=base, font=\small] at (\x, -1.27) {\x};
  \node[font=\small, black!65] at (5.5, -1.5) {member};
  #1
\end{tikzpicture}}



% An answer, set apart from the question it belongs to.
\newcommand{\ans}[1]{\textbf{\textcolor{pltblue!75!black}{#1}}}
\newcommand{\arrow}{\textcolor{pltblue!75!black}{$\Rightarrow$}\;}
\newcommand{\mans}[1]{{\color{pltblue!75!black}#1}}

\begin{document}

\begin{center}
{\ppfont\LARGE Answers}
\end{center}

{\small The sheet with its answers, \ans{in blue}. Parts 1 and 2 are done
before the lecture, Part 3 after it.}

\parthead{1}{A walker with no map}

Four friends: A, B, C and D. A walker starts at A and at every step moves to
one friend of the person it is standing on, each with the same chance.

\begin{center}
\resizebox{0.5\textwidth}{!}{\fournet{%
  \node[text=pltblue!75!black, font=\bfseries] at (bA) {3/8};
  \node[text=pltblue!75!black, font=\bfseries] at (bB) {2/8};
  \node[text=pltblue!75!black, font=\bfseries] at (bC) {2/8};
  \node[text=pltblue!75!black, font=\bfseries] at (bD) {1/8};}}\\[0.2em]
{\small the boxes are filled in Question 6}
\end{center}

{\bf Question 1}: Before you compute anything, guess. After a very long walk,
at which of the four people is the walker found most often? Why?

\arrow Any guess is fine; Question 6 checks it. The answer is \ans{A}, the
person with the most friends. A common wrong guess is D, because every walk
that reaches D has to come straight back through A.

{\bf Question 2}: In each cell of the table, write the chance that the walker
moves from the person of that row to the person of that column in one step.
Then add up each row.

\begin{center}
{\setlength{\tabcolsep}{6pt}
\begin{tabular}{|C{2.4cm}|C{1.5cm}|C{1.5cm}|C{1.5cm}|C{1.5cm}|C{1.5cm}|}
\hline
\small from \textbackslash{} to & \bfseries A & \bfseries B & \bfseries C & \bfseries D & \small sum \\
\hline
\bfseries A & \mans{0} & \mans{$\frac13$} & \mans{$\frac13$} & \mans{$\frac13$} & \mans{1} \\[0.12cm] \hline
\bfseries B & \mans{$\frac12$} & \mans{0} & \mans{$\frac12$} & \mans{0} & \mans{1} \\[0.12cm] \hline
\bfseries C & \mans{$\frac12$} & \mans{$\frac12$} & \mans{0} & \mans{0} & \mans{1} \\[0.12cm] \hline
\bfseries D & \mans{1} & \mans{0} & \mans{0} & \mans{0} & \mans{1} \\[0.12cm] \hline
\end{tabular}}
\end{center}

\arrow Each row is one over the number of friends of that row's person, in
the columns of those friends. \ans{Every row adds up to 1}: the walker has to
go somewhere. This table is the matrix $P$ of the lecture (slide 14), one row
per starting person.

{\bf Question 3}: The walker starts at A. What is the chance that it is at B
after exactly three steps?

\begin{center}
{\setlength{\tabcolsep}{8pt}
\begin{tabular}{|C{5.0cm}|C{5.0cm}|}
\hline
\bfseries route & \bfseries chance \\ \hline
A\,$\to$\,D\,$\to$\,A\,$\to$\,B & $\frac13 \times 1 \times \frac13 = \frac19$ \\[0.08cm] \hline
\mans{A\,$\to$\,B\,$\to$\,A\,$\to$\,B} & \mans{$\frac13 \times \frac12 \times \frac13 = \frac1{18}$} \\[0.08cm] \hline
\mans{A\,$\to$\,C\,$\to$\,A\,$\to$\,B} & \mans{$\frac13 \times \frac12 \times \frac13 = \frac1{18}$} \\[0.08cm] \hline
\mans{A\,$\to$\,B\,$\to$\,C\,$\to$\,B} & \mans{$\frac13 \times \frac12 \times \frac12 = \frac1{12}$} \\[0.08cm] \hline
& \\[0.3cm] \hline
\end{tabular}}
\end{center}

\arrow Four routes, and the last row stays empty.
$\frac19 + \frac1{18} + \frac1{18} + \frac1{12} = \frac{4+2+2+3}{36} =$
\ans{$\frac{11}{36} \approx 0.31$}. The routes that go back to A are the ones
people miss: the walker has no memory, so it may undo its own step.

\parthead{2}{Drunk at the karate club}

A university karate club in the 1970s, with 34 members and 78 friendships. The
simulator (\simurl) runs the same walker on it.

{\bf Question 4}: Press \textbf{Run 5,000 steps} four times: 20,000 steps.
Write the most and the least visited node in the table, with its edges and its
share. Then divide the edges by the share.

\begin{center}
{\setlength{\tabcolsep}{6pt}
\begin{tabular}{|C{3.2cm}|C{1.7cm}|C{1.7cm}|C{2.0cm}|C{2.6cm}|}
\hline
& \bfseries node & \bfseries edges & \bfseries share & \bfseries edges $\div$ share \\ \hline
\bfseries most visits & \mans{33} & \mans{17} & \mans{0.109} & \mans{156} \\[0.08cm] \hline
\bfseries fewest visits & \mans{11} & \mans{1} & \mans{0.0064} & \mans{156} \\[0.08cm] \hline
\end{tabular}}
\end{center}

\arrow These are the long-run values. One run of 20,000 steps gives numbers
near them, not equal to them: the ratio for node 33 usually lands between
about 145 and 165, and for node 11 it can be 10\% off. The most visited node
is node 33 in about three runs out of four and node 0 (16 edges, share 0.103)
in the others. Both give 156.

{\bf Question 5}: Are the two numbers in the last column the same? How are
they related to the 78 friendships?

\arrow \ans{Nearly the same, and the same for everybody in the room: about 156 =
2 $\times$ 78.} The gaps between neighbours are the luck of one run, and they
shrink if you run longer. So the share of time at a node is its number of
friends divided by twice the number of friendships:
\[
\text{share} = \frac{\text{edges}}{2 \times \text{friendships}} = \frac{k}{2m}.
\]
It is twice, not once, because every friendship has two ends, and the walker
can arrive through either one.

{\bf Question 6}: The four friends in Part 1 have 4 friendships. Work out the
share of time at A, B, C and D. Was your guess in Question 1 right?

\arrow $2m = 8$: \ans{A $= \frac38$, B $= \frac28$, C $= \frac28$, D $=
\frac18$} (the boxes on the drawing at the top). A gets most because it has the
most friends. D gets least even though every visit to D is followed by a visit
to A: D has one friend, and one friend is all that counts.

\parthead{3}{Two rooms}

Twelve members and 30 friendships. The notebook (\molaburl) runs the cell
printed on the sheet and prints
\begin{center}
{\small\begin{tabular}{@{}l@{\;}l@{}}
\texttt{eigenvalues:} & \texttt{[1, 0.889, 0, 0, 0, -0.088, -0.2, -0.2, -0.266, -0.333, -0.379, -0.423]}\\
\texttt{second eigenvector:} & \texttt{[0.98, 1, 0.74, 0.98, 1, 0.74, -1, -0.98, -0.74, -1, -0.98, -0.74]}
\end{tabular}}
\end{center}
(shortened here). On some machines every sign of the second eigenvector is the
other way round. That is fine: an eigenvector times $-1$ is still an
eigenvector, and the answers below only swap the two groups.

{\bf Question 7}: Mark the twelve numbers of the second eigenvector as dots on
the axes.

\begin{center}
\resizebox{0.78\textwidth}{!}{\eigaxes{%
  \foreach \x/\y in {0/0.98, 1/1, 2/0.74, 3/0.98, 4/1, 5/0.74,
                     6/-1, 7/-0.98, 8/-0.74, 9/-1, 10/-0.98, 11/-0.74}
    \node[dot] at (\x, \y) {};}}
\end{center}

{\bf Question 8}: On the drawing, darken every member whose number is above
zero. What do the darkened members have in common?

\begin{center}
\resizebox{0.62\textwidth}{!}{\tworooms{fill=pltblue!45}{}}
\end{center}

\arrow \ans{Members 0 to 5: one of the two groups.} Inside each group almost
everybody is friends with everybody, and only two friendships cross between
them. The sign of the second eigenvector names the group, which is why it is
called a community signal (slide 44).

{\bf Question 9}: Which four members have the numbers closest to zero? Where
are they on the drawing?

\arrow \ans{Members 2, 5, 8 and 11 ($\pm$0.74): the ends of the two friendships
that cross between the groups.} A walker at one of them is a single step from
the other group, so it belongs to its own group a little less firmly.

{\bf Question 10}: Work out 0.889\textsuperscript{5}, 0.889\textsuperscript{40}
and 0.423\textsuperscript{5}. After 5 steps, which eigenvalues still have a
term in the walk? And after 40 steps?

\arrow \ans{0.889\textsuperscript{5} $\approx$ 0.56,
0.889\textsuperscript{40} $\approx$ 0.009, 0.423\textsuperscript{5} $\approx$
0.013.} Every other eigenvalue is smaller in size than 0.423, so its term is
even smaller. \ans{After 5 steps two terms are left: the eigenvalue 1, which
never shrinks, and 0.889. After 40 steps only the eigenvalue 1 is left.}

{\bf Question 11}: The lecture asked: short walks show communities, long walks
show degree. Why? Answer in two sentences.

\arrow \ans{Long walk: every term except the eigenvalue 1 has died, and that
term is the share of time from Part 2, edges $\div$ 2m, so all that is left is
degree. Short walk: the second term, 0.889\textsuperscript{t}, is still large
for the first ten or so steps, and its eigenvector has one sign in each group
(Part 3), so the walker still remembers which group it started in.} The ten
steps are the relaxation time of the lecture, $\tau = 1/(1-0.889) = 9$.

{\bf Question 12} (finished early): Delete the friendship 2--8 and run the cell
again. Is the second eigenvalue closer to 1, or further from 1? Does a walker
that starts at member 0 stay among members 0 to 5 for longer?

\arrow \ans{Closer: 0.946. Longer.} The term 0.946\textsuperscript{t} shrinks
more slowly than 0.889\textsuperscript{t} (0.76 against 0.56 after 5 steps),
so the walker remembers its group for longer: with one way across instead of
two, it takes longer to find the way out. The relaxation time grows from 9 to
$1/(1-0.946) \approx 19$ steps. This is the bottleneck of slides 34--37.

\end{document}
